\chapter{النوم واليقظة}
\label{chap:sleep}

تستدعي العمليات بداخل النواة آليات \indextext{النوم واليقظة (Sleep & Wakeup)} عندما تحتاج إلى الانتظار حتى وقوع حدث معين (مثل قراءة بيانات من القرص أو انتهاء عملية ابن).

تُوفر النواة دالتين أساسيتين لإدارة الانتظار التزامني:
\begin{itemize}
\item \texttt{sleep(chan, lock)}: تغير وضع العملية الحالية إلى \texttt{SLEEPING} وتسجل عنوان قناة الانتظار \texttt{chan}، وتُحرر القفل الممرر \texttt{lock}، ثم تستدعي المجدول لنقل المعالج لعملية أخرى.
\item \texttt{wakeup(chan)}: تبحث بداخل جدول العمليات عن جميع العمليات النائمة بداخل نفس قناة الانتظار \texttt{chan}، وتغير وضعيتها إلى \texttt{RUNNABLE} لتصبح جاهزة للتنفيذ.
\end{itemize}

تعتمد آليات النوم واليقظة في \texttt{xv6} قفل جدول العمليات \texttt{p->lock} لتجنب المشكلة التاريخية المشهورة بـ \indextext{الاستيقاظ الضائع (Lost Wakeup)}؛
حيث يضمن القفل غياب الفجوة الزمنية بين فحص شرط الانتظار والدخول الفعلي في حالة النوم.
